\documentclass[11pt]{article}

\usepackage[margin=1in]{geometry}
\usepackage{amsmath,amssymb,amsthm}
\usepackage{graphicx}
\usepackage[hidelinks]{hyperref}

\newtheorem{theorem}{Theorem}
\newtheorem{lemma}{Lemma}

\newcommand{\N}{\mathbb N}
\newcommand{\R}{\mathbb R}
\newcommand{\dist}{\operatorname{dist}}
\newcommand{\diam}{\operatorname{diam}}
\newcommand{\spanop}{\operatorname{span}}

\title{Zero Thickness Does Not Force a Hilbert Bi-Lipschitz Embedding}
\author{FA/Banach run fa\_banach\_001}
\date{June 14, 2026}

\begin{document}
\maketitle

\section*{Status}

This packet gives a counterexample to the second open question of
Margaris--Robinson \cite{MR2018}. Thickness exponent zero does not force the
existence of a nonlinear bi-Lipschitz embedding into a Hilbert space. Since
Euclidean spaces are Hilbert spaces, it also rules out the Euclidean target.

\section*{Source Question}

After proving an embedding theorem for compact subsets of Banach spaces with
thickness exponent less than \(1\), the source paper asks:

\begin{figure}[h]
\centering
\includegraphics[width=\textwidth]{figures/open_problem_crop.png}
\caption{Source evidence from arXiv:1811.03872, PDF page 20.}
\end{figure}

The construction below addresses the second question.

\section*{The Counterexample}

Let \(Q_m=\{0,1\}^m\) be the \(m\)-dimensional Hamming cube, equipped with the
Hamming metric \(d_m\). Let \(0_m\) denote its zero vertex. Put
\[
  E_m=\ell_\infty(Q_m)
\]
and embed \(Q_m\) into \(E_m\) by the pointed Kuratowski map
\[
  \iota_m(x)=\big(d_m(x,y)-d_m(0_m,y)\big)_{y\in Q_m}.
\]
This map is isometric: for every \(x,x'\in Q_m\),
\[
  \|\iota_m(x)-\iota_m(x')\|_\infty
  =\sup_{y\in Q_m}|d_m(x,y)-d_m(x',y)|
  =d_m(x,x'),
\]
where the reverse inequality is the triangle inequality and equality follows
by taking \(y=x\).

Let
\[
  B=\left(\bigoplus_{m=1}^\infty E_m\right)_{c_0}
\]
with norm \(\|(u_m)\|=\sup_m\|u_m\|_{E_m}\). Choose
\[
  R_m=2^{-m^3}, \qquad s_m=\frac{R_m}{m}.
\]
For \(x\in Q_m\), let \(u_{m,x}\in B\) be the vector whose \(m\)-th coordinate
is \(s_m\iota_m(x)\) and whose other coordinates are zero. Define
\[
  X=\{0\}\cup \{u_{m,x}:m\in\N,\ x\in Q_m\}.
\]

\begin{theorem}
The compact set \(X\subset B\) has thickness exponent \(0\), but \(X\) does
not admit a bi-Lipschitz embedding into any Hilbert space.
\end{theorem}

\begin{proof}
First, \(X\) is compact. Each level \(\{u_{m,x}:x\in Q_m\}\) is finite, and
\[
  \|u_{m,x}\|\le s_m\diam(Q_m)=s_m m=R_m\to 0.
\]
Thus all levels eventually lie in an arbitrarily small ball around \(0\).

We next compute a box-counting bound. If
\[
  R_{M+1}<\varepsilon\le R_M,
\]
then all levels \(m>M\) are contained in the \(\varepsilon\)-ball around \(0\).
The first \(M\) levels can be covered point by point. Hence
\[
  N(X,\varepsilon)\le 1+\sum_{m=1}^M |Q_m|
  =1+\sum_{m=1}^M 2^m
  \le 2^{M+1}.
\]
Since \(-\log R_M=M^3\log 2\), it follows that
\[
  \overline{\dim}_B(X)
  =\limsup_{\varepsilon\downarrow 0}
      \frac{\log N(X,\varepsilon)}{-\log\varepsilon}
  \le
  \limsup_{M\to\infty}\frac{(M+1)\log 2}{M^3\log 2}=0.
\]
The thickness exponent is always bounded above by the upper box-counting
dimension: an \(\varepsilon\)-cover by \(N(X,\varepsilon)\) balls gives an
\(\varepsilon\)-approximating subspace spanned by the \(N(X,\varepsilon)\)
centres. Therefore \(\tau(X,B)=0\).

It remains to exclude Hilbert embeddings. Suppose, towards a contradiction,
that \(F:X\to H\) is bi-Lipschitz into a Hilbert space \(H\). Then there are
constants \(0<a\le b<\infty\) such that
\[
  a\|u-v\|_B\le \|F(u)-F(v)\|_H\le b\|u-v\|_B
  \qquad (u,v\in X).
\]
Restrict \(F\) to the \(m\)-th level. Since \(\iota_m\) is isometric, for
all \(x,x'\in Q_m\),
\[
  a s_m d_m(x,x')
  \le \|F(u_{m,x})-F(u_{m,x'})\|_H
  \le b s_m d_m(x,x').
\]
After dividing by \(s_m\), this gives a Hilbert embedding of \(Q_m\) with
distortion at most \(b/a\), uniformly in \(m\).

The following standard Enflo cube estimate contradicts this. If
\(G:Q_m\to H\) has Lipschitz constant at most \(\beta\) and inverse
Lipschitz constant at least \(\alpha\), then \(\beta/\alpha\ge \sqrt m\).
Indeed, identify \(Q_m\) with \(\{-1,1\}^m\). For \(\varepsilon\in Q_m\), let
\(-\varepsilon\) be the antipodal vertex and \(\varepsilon^{(j)}\) the vertex
obtained by flipping the \(j\)-th coordinate. Hilbert-valued Walsh expansion
gives
\[
  \mathbb E_\varepsilon \|G(\varepsilon)-G(-\varepsilon)\|^2
  \le
  \sum_{j=1}^m
  \mathbb E_\varepsilon \|G(\varepsilon)-G(\varepsilon^{(j)})\|^2 .
\]
The right side is at most \(m\beta^2\), since each edge has Hamming length
\(1\). Therefore some antipodal pair satisfies
\[
  \|G(\varepsilon)-G(-\varepsilon)\|\le \beta\sqrt m.
\]
But the lower Lipschitz bound gives
\[
  \|G(\varepsilon)-G(-\varepsilon)\|\ge \alpha m.
\]
Thus \(\beta/\alpha\ge \sqrt m\).

Applying this to the uniformly distorted embeddings of the cubes \(Q_m\)
forces \(b/a\ge \sqrt m\) for all \(m\), impossible. Hence no such Hilbert
embedding \(F\) exists.
\end{proof}

\section*{Verification Notes}

The construction deliberately separates global and local geometry. The levels
shrink quickly enough that the global upper box-counting dimension is zero,
but each level retains the full Hamming-cube metric after rescaling. Thickness
sees the global covering rate; Hilbert bi-Lipschitz embeddability would have
to handle all finite cube levels with one distortion constant.

\section*{Novelty and Search Bounds}

Before packaging, the run's cheap indexes were searched for arXiv:1811.03872,
``zero thickness'', ``bi-Lipschitz'', ``Euclidean'', ``Hilbert'', and
``Assouad''. No prior packet or attempt for this target was found. The proof
uses the standard Enflo lower bound for Hamming cubes in Hilbert space and is
packaged as an original counterexample to the source question.

\section*{Human Review Recommendation}

Review as a counterexample packet. The essential checks are the
box-counting/thickness computation and the uniform-distortion obstruction from
the Hamming cube levels.

\begin{thebibliography}{9}

\bibitem{MR2018}
A. Margaris and J. C. Robinson,
\emph{Embedding properties of sets with finite box-counting dimension},
arXiv:1811.03872.

\end{thebibliography}

\end{document}
